\documentclass[11pt]{article}
\usepackage{amsmath,amssymb}

\begin{document}

\section*{Line Graphs of Trees}

We searched for trees \(T\) such that the chromatic symmetric function of the
line graph \(L(T)\) fails to be \(e\)-positive. Using the tree data in
\texttt{trees\(n\).g6}, the first counterexample occurs on \(7\) vertices.

More precisely, consider the tree with edge set
\[
\{\{0,3\},\{1,4\},\{2,5\},\{0,6\},\{1,6\},\{2,6\}\}.
\]
This is the \(3\)-legged spider with center \(6\) and each leg of length \(2\).
Its line graph has edge set
\[
\{\{0,3\},\{1,4\},\{2,5\},\{3,4\},\{3,5\},\{4,5\}\}.
\]

For this graph,
\[
X_{L(T)}
=
6e_{3,2,1}
-6e_{3,3}
+6e_{4,1,1}
+12e_{4,2}
+18e_{5,1}
+12e_{6}.
\]
In particular,
\[
[e_{3,3}]\,X_{L(T)}=-6,
\]
so \(X_{L(T)}\) is not \(e\)-positive.

Its Schur expansion is
\[
X_{L(T)}
=
48s_{1,1,1,1,1,1}
+48s_{2,1,1,1,1}
+24s_{2,2,1,1}
+12s_{3,1,1,1}
+6s_{3,2,1}.
\]

\bigskip

It is also natural to restrict to trees in which every non-leaf vertex has
degree \(3\). In this family, the first counterexample occurs on \(10\)
vertices. Consider the tree with edge set
\[
\{\{0,6\},\{1,6\},\{2,7\},\{3,7\},\{4,8\},\{5,8\},\{6,9\},\{7,9\},\{8,9\}\}.
\]
This is the cubic tree with four internal vertices \(6,7,8,9\), where the
vertex \(9\) is adjacent to \(6,7,8\), and each of \(6,7,8\) is adjacent to two
leaves.

Its line graph has edge set
\[
\{\{0,1\},\{0,6\},\{1,6\},\{2,3\},\{2,7\},\{3,7\},\{4,5\},\{4,8\},\{5,8\},
\{6,7\},\{6,8\},\{7,8\}\}.
\]
For this graph,
\[
\begin{aligned}
X_{L(T)}={}&48e_{3,3,3}+96e_{4,3,2}-48e_{4,4,1}+48e_{5,2,2}+192e_{5,3,1}-144e_{5,4}\\
&+144e_{6,2,1}+288e_{6,3}+96e_{7,1,1}+192e_{7,2}+240e_{8,1}+144e_9.
\end{aligned}
\]
In particular,
\[
[e_{4,4,1}]\,X_{L(T)}=-48,
\qquad
[e_{5,4}]\,X_{L(T)}=-144,
\]
so \(X_{L(T)}\) is not \(e\)-positive.

Its Schur expansion is
\[
\begin{aligned}
X_{L(T)}={}&1296s_{1,1,1,1,1,1,1,1,1}
+1728s_{2,1,1,1,1,1,1,1}
+1584s_{2,2,1,1,1,1,1}
+1152s_{2,2,2,1,1,1}\\
&+288s_{2,2,2,2,1}
+576s_{3,1,1,1,1,1,1}
+672s_{3,2,1,1,1,1}
+528s_{3,2,2,1,1}
+96s_{3,2,2,2}\\
&+192s_{3,3,1,1,1}
+192s_{3,3,2,1}
+48s_{3,3,3}.
\end{aligned}
\]

\bigskip

We also computed the chromatic symmetric functions of the line graphs
\(L(K_{2,n})\) for small \(n\). In the cases below, all coefficients in the
\(e\)-basis are positive. The last column records the sum of the coefficients in
the \(e\)-expansion.

\medskip

\small
\begin{center}
\begin{tabular}{c|p{10cm}|c}
\(n\) & \(X_{L(K_{2,n})}\) in the \(e\)-basis & sum of \(e\)-coefficients \\
\hline
1 &
\(2e_2\) &
\(2\) \\
2 &
\(2e_{2,2}+12e_4\) &
\(14\) \\
3 &
\(12e_{3,3}+12e_{4,2}+24e_{5,1}+156e_6\) &
\(204\) \\
4 &
\(216e_{4,4}+96e_{5,3}+336e_{6,2}+864e_{7,1}+3504e_8\) &
\(5016\) \\
5 &
\(5280e_{5,5}+2160e_{6,4}+5760e_{7,3}+14640e_{8,2}+37440e_{9,1}+120240e_{10}\) &
\(185520\)
\end{tabular}
\end{center}
\normalsize

\bigskip

We also computed the chromatic symmetric functions of the line graphs of the
\(2\times n\) ladder graphs \(P_2 \square P_n\) for small \(n\). Again, in the
cases below all coefficients in the \(e\)-basis are positive.

\medskip

\small
\begin{center}
\begin{tabular}{c|p{10cm}|c}
\(n\) & \(X_{L(P_2 \square P_n)}\) in the \(e\)-basis & sum of \(e\)-coefficients \\
\hline
2 &
\(2e_{2,2}+12e_4\) &
\(14\) \\
3 &
\(2e_{3,2,2}+4e_{3,3,1}+6e_{4,2,1}+26e_{4,3}+82e_{5,2}+80e_{6,1}+238e_7\) &
\(438\) \\
4 &
\(\begin{aligned}
&6e_{3,3,2,2}
+26e_{4,3,2,1}
+54e_{4,3,3}
+18e_{4,4,1,1}
+64e_{4,4,2}\\
&+18e_{5,2,2,1}
+278e_{5,3,2}
+264e_{5,4,1}
+730e_{5,5}
+10e_{6,2,1,1}\\
&+256e_{6,2,2}
+258e_{6,3,1}
+768e_{6,4}
+686e_{7,2,1}
+990e_{7,3}\\
&+252e_{8,1,1}
+2432e_{8,2}
+2748e_{9,1}
+3840e_{10}
\end{aligned}\) &
\(13698\)
\end{tabular}
\end{center}
\normalsize

\bigskip

Finally, we considered the following Ferrers-diagram family. For a partition
\(\lambda\), let \(G_\lambda\) be the graph whose vertices are the cells of the
Ferrers diagram of \(\lambda\), with two cells adjacent if they share an edge.
We then computed \(X_{L(G_\lambda)}\) in the \(e\)-basis. For all Ferrers
shapes with at most \(7\) cells, the result was \(e\)-positive.

\medskip

\small
\begin{center}
\begin{tabular}{c|p{9.5cm}|c}
number of cells & Ferrers shapes \(\lambda\) checked & result for \(X_{L(G_\lambda)}\) \\
\hline
1 & \((1)\) & all \(e\)-positive \\
2 & \((2)\), \((1,1)\) & all \(e\)-positive \\
3 & \((3)\), \((2,1)\), \((1,1,1)\) & all \(e\)-positive \\
4 & \((4)\), \((3,1)\), \((2,2)\), \((2,1,1)\), \((1,1,1,1)\) & all \(e\)-positive \\
5 & \((5)\), \((4,1)\), \((3,2)\), \((3,1,1)\), \((2,2,1)\), \((2,1,1,1)\), \((1,1,1,1,1)\) & all \(e\)-positive \\
6 & \((6)\), \((5,1)\), \((4,2)\), \((4,1,1)\), \((3,3)\), \((3,2,1)\), \((3,1,1,1)\), \((2,2,2)\), \((2,2,1,1)\), \((2,1,1,1,1)\), \((1,1,1,1,1,1)\) & all \(e\)-positive \\
7 & \((7)\), \((6,1)\), \((5,2)\), \((5,1,1)\), \((4,3)\), \((4,2,1)\), \((4,1,1,1)\), \((3,3,1)\), \((3,2,2)\), \((3,2,1,1)\), \((3,1,1,1,1)\), \((2,2,2,1)\), \((2,2,1,1,1)\), \((2,1,1,1,1,1)\), \((1,1,1,1,1,1,1)\) & all \(e\)-positive
\end{tabular}
\end{center}
\normalsize

\end{document}
